Most modern neural networks rely on deterministic activation functions like ReLU, which lack inherent stochasticity for regularization or uncertainty estimation. While stochastic methods like Dropout or noise-based activations exist, the potential of using categorical sampling at intermediate layers remains largely unexplored. In this work, we propose Softmax Sampling Activation (SSA), a technique that treats intermediate activations as logits and samples a neuron to pass forward according to its relative probability. We leverage the Gumbel-Softmax reparameterization trick to enable differentiable categorical sampling within the network architecture. Our experiments on CIFAR-10 using a VGG-11 architecture reveal that while applying SSA at every layer leads to vanishing signals and training instability, selective application at the final feature layer allows for effective training. We show that "soft" SSA (continuous relaxation) achieves a test accuracy of 83.09\%, compared to 86.46\% for a standard ReLU baseline, whereas "hard" SSA (discrete sampling) fails to converge. Our analysis suggests that the competitive nature of categorical sampling creates a bottleneck that, while regularizing, can restrict representational capacity. We conclude that SSA is a viable but challenging primitive that provides a natural mechanism for stochastic feature selection.
